\section{Task Definition and Data}\label{sec:task}
\paragraph{Task.} At prediction time $t$ (every 15 minutes) we predict $y$, the traffic-weighted
mean car travel time (min) on the target segment over $[T,T{+}15\text{ min})$, $T=t+60$ min. With
$\mathit{TT}_i$ the median travel time (s) and $n_i$ the vehicle count of 5-minute bin $i$,
$y=\sum_i \mathit{TT}_i n_i / \sum_i n_i / 60$, the time a typical car experiences. One row is one
(day, $t$) pair. The horizon and step match the decision ``leave now or 30--60 minutes later''.
We keep $T\in[06{:}00,21{:}00)$: night traffic is free-flowing, would inflate accuracy, and lies
outside the decision window.

\paragraph{Data.} Travel times come from the Freeway Bureau TDCS \textsc{M04A} table (median
gantry-to-gantry travel time per vehicle type, 5-minute bins), passenger cars only (commuters
drive cars). The target segment is \res{cfg.seg.target} (88.0K--92.8K); the upstream segment is
\res{cfg.seg.upstream} (75.0K--88.0K), immediately to the north. Hourly rainfall comes from the CWA
Hsinchu station and working days from the DGPA calendar. We verified on raw files that travel time
is in seconds and that time stamps mark the \emph{start} of each bin. Training covers
\res{cfg.trainstart} to \res{cfg.trainend} (\res{data.daystrain} days, \res{data.rowstrain} rows,
two Lunar New Years and a typhoon season); the held-out test set, the most recent data as in
deployment, is \res{cfg.teststart} to \res{cfg.testend} (\res{data.daystest} days,
\res{data.rowstest} rows) and was evaluated once, after the design was frozen. No 5-minute file
was missing and \res{data.validPct}\% of bins are valid (travel time and count $>0$), so no row was
dropped or imputed. $y$ ranges from \res{y.min} to \res{y.max}\,min (median \res{y.median}, 95th
percentile \res{y.p95}; Table~\ref{tab:summary}). We keep the extremes: they are real jams, the
cases a commuter most wants to foresee.

\begin{table}[t]\centering\small
\fitinput{tables/tab_data_summary}
\caption{Summary statistics of the target and key features (training rows).}\label{tab:summary}
\end{table}
